\chapter*{Hướng Dẫn Sử Dụng Sách}
\addcontentsline{toc}{chapter}{Hướng Dẫn Sử Dụng Sách}
\markboth{Hướng Dẫn Sử Dụng Sách}{Hướng Dẫn Sử Dụng Sách}

\begin{truyen}{Cách Đọc Cuốn Này}{Read by pair, reflect by contrast.}
\chuhoa{Đ}{ừng đọc kiểu lướt cho hết trang. Cuốn này mà đọc như nuốt tin ngắn thì chỉ còn mấy câu nghe hay tai. Mỗi chương là một cặp kéo qua kéo lại, nên phải đọc chậm, ngẫm, rồi đem so với chính đời mình thì mới ra chuyện.}
\textit{(This book works best when read slowly. Each chapter is a symmetrical pair, and each pair should be held side by side to reveal its meaning. Reading quickly may give the main idea, but reading slowly against your own life makes it deeper.)}

\textbf{Người dạy học:} có thể lấy từng tiểu mục làm mồi để gợi thảo luận, viết đoạn văn, hoặc nói chuyện chủ nhiệm cho bớt nhạt.
\textit{(\textbf{Teachers:} may use each mini-section as a short scenario for discussion, reflective writing, or homeroom activities.)}

\textbf{Người đọc cá nhân:} tốt nhất đọc 1--2 tiểu mục một ngày, gạch một câu thấy nhói, rồi tự viết vài dòng thật về mình.
\textit{(\textbf{Individual readers:} should read 1--2 mini-sections a day, mark one sentence, then write down one honest response.)}
\end{truyen}

\begin{baihoc}
Cuốn này không viết để trưng triết lý. Nó viết để người đọc đụng đâu thấy mình ở đó.
\textit{(This book is not for finishing quickly, but for pausing at the right moment. Each symmetrical pair is a mirror, not an exam.)}
\end{baihoc}

\begin{ghinhoanh}
\textbf{Reading tips:}

Read by contrast, not by hurry.

A short honest pause is better than a long careless reading.
\end{ghinhoanh}

\ngancach
